\documentclass[11pt,a4paper]{article}
\usepackage[T1]{fontenc}
\usepackage[utf8]{inputenc}
\usepackage{lmodern,amsmath,amssymb}
\usepackage[margin=25mm]{geometry}
\usepackage[hidelinks]{hyperref}
\hypersetup{pdftitle={A physical cut: elastic reversal, finite speed and memory},pdfauthor={navstokgap research working notes}}
\newcommand{\tightlist}{\setlength{\itemsep}{0pt}\setlength{\parskip}{0pt}}
\setlength{\emergencystretch}{3em}
\setcounter{secnumdepth}{0}
\begin{document}
\hypertarget{a-physical-cut-elastic-reversal-finite-speed-and-memory}{%
\section{A physical cut: elastic reversal, finite speed and
memory}\label{a-physical-cut-elastic-reversal-finite-speed-and-memory}}

A05, reviewed with B13, 2026-09-07. An equal-mass elastic collision
realizes a random midpoint with exact energy and momentum conservation.
Its action cost tends to zero with the interval duration at fixed speed.
A sharp midpoint bound extends this conclusion to every bounded-speed
bridge law. Finally, a position-only convolution law with ballistic
support is necessarily a deterministic drift. Together these tests
identify velocity memory as a concrete ingredient for the next
refinement model.

\hypertarget{one-collision-with-a-momentum-receiver}{%
\subsection{1. One collision with a momentum
receiver}\label{one-collision-with-a-momentum-receiver}}

Fix \(m>0\), \(0<u<c\) and duration \(\Delta>0\). Choose a common random
sign \(\sigma=\pm1\) with equal probabilities. At time zero, a labelled
tracer is at \(0\) with velocity \(\sigma u\), and a second particle of
the same mass is at \(\sigma u\Delta\) with velocity \(-\sigma u\). Both
move freely on the line until their instantaneous elastic collision at
time \(\Delta/2\). The equal-mass collision exchanges their velocities.
Thus the tracer follows

\[X(t)=\begin{cases}\sigma ut,&0\le t\le\Delta/2,\\
\sigma u(\Delta-t),&\Delta/2\le t\le\Delta.\end{cases}\]

Its midpoint \(Y=\sigma u\Delta/2\) is random, while
\(X(0)=X(\Delta)=0\). The other particle returns to its own initial
position. Throughout the motion, total momentum is zero and total
kinetic energy is \(mu^2\). Each labelled velocity changes sign at the
collision, with opposite impulses of magnitude \(2mu\). Both particles
stay below \(c\) in the chosen frame.

The tracer's free kinetic action relative to the stationary path with
the same position endpoints is

\[D=\frac m2\int_0^\Delta\dot X^2dt=\frac{mu^2\Delta}{2}.
\qquad \operatorname{Var}Y=\frac{u^2\Delta^2}{4}.\]

Define the midpoint action parameter in A03's normalization by
\(\kappa_{\rm mid}=4m\operatorname{Var}Y/\Delta\). This gives
\(\kappa_{\rm mid}=mu^2\Delta\) and \(\mathbb E D=\kappa_{\rm mid}/2\).
At fixed energy and speed, both tend to zero as \(\Delta\downarrow0\).
Keeping a supplied \(\kappa_{\rm mid}>0\) instead requires
\(u^2=\kappa_{\rm mid}/(m\Delta)\) and pair energy
\(\kappa_{\rm mid}/\Delta\).

This is an exactly specified hard-collision model with random initial
data. The timing follows from the prepared initial separation
\(u\Delta\), rather than a new stochastic clock. Changing \(\Delta\)
changes that separation. The stationary comparison shares positions and
times, not initial velocities or apparatus state. Consequently the
construction preserves the coarse position endpoints, but does not
implement an invisible intervention on a fixed full phase-space
preparation. Smooth collision potentials and a Lorentz-covariant
dynamics are separate models.

\hypertarget{sharp-midpoint-support-and-variance-bound}{%
\subsection{2. Sharp midpoint support and variance
bound}\label{sharp-midpoint-support-and-variance-bound}}

Let any absolutely continuous path on \([0,\Delta]\) have endpoints
\(x,z\) and \(|\dot X|\le u<c\) almost everywhere. Write
\(v=(z-x)/\Delta\), with \(|v|\le u\), and
\(\zeta=X(\Delta/2)-(x+z)/2\). Each half-interval gives

\[X(\Delta/2)\in[x-u\Delta/2,x+u\Delta/2]
\cap[z-u\Delta/2,z+u\Delta/2].\]

This intersection is centered at \((x+z)/2\) with half-width
\(r=\Delta(u-|v|)/2\). Therefore any probability law on these admissible
paths satisfies

\[\operatorname{Var}Y\le r^2,
\qquad 0\le\kappa_{\rm mid}:=\frac{4m}{\Delta}\operatorname{Var}Y
\le m\Delta(u-|v|)^2.\]

Proof of the variance bound:
\(\operatorname{Var}Y\le \mathbb E(Y-(x+z)/2)^2\le r^2\). Equal
probabilities at the two extremal midpoints saturate it; the
corresponding two-segment paths obey the speed bound. For \(x=z\) the
elastic example realizes this sharp case.

The two-segment interpolant through \(Y\) has excess kinetic action
\(D_{\rm poly}=2m\zeta^2/\Delta\). When the midpoint law is centered on
\((x+z)/2\), \(\mathbb E D_{\rm poly}=\kappa_{\rm mid}/2\). For a biased
midpoint there is the additional term \(2m(\mathbb E\zeta)^2/\Delta\).
This distinction avoids identifying a variance with all action costs.

A Gaussian midpoint with variance \(\kappa\Delta/(4m)\) has unbounded
support for every \(\kappa>0\) and so fails the exact speed constraint
at every duration. Even a bounded replacement matching only that
variance requires \(\kappa\le m\Delta(u-|v|)^2\). The latter inequality
gives a necessary resolution restriction, not a derivation of
\(\kappa\).

At \(|v|=u\) the allowable midpoint is unique. Indeed a path achieving
the maximum displacement \(u\Delta\) must have velocity \(u\) almost
everywhere. Thus after choosing an extremal midpoint in the return-path
example, each half is already a saturated straight segment: finer
admissible sampling inside that half is deterministic. Reapplying a
nonzero midpoint fluctuation rule independently at every cut would
change the coarse law or violate speed.

\hypertarget{finite-speed-and-a-position-only-convolution-law}{%
\subsection{3. Finite speed and a position-only convolution
law}\label{finite-speed-and-a-position-only-convolution-law}}

\textbf{Proposition.} Suppose probability measures \(\mu_t\) on the line
satisfy \(\mu_0=\delta_0\), \(\mu_{s+t}=\mu_s*\mu_t\), and
\(\operatorname{supp}\mu_t\subset[-ut,ut]\) for every \(t\ge0\), with a
fixed finite \(u\). Then \(\mu_t=\delta_{bt}\) for some \(|b|\le u\).

\textbf{Proof.} For fixed \(T>0\) and every positive integer \(n\),
convolution and bounded support give

\[\operatorname{Var}\mu_T=n\operatorname{Var}\mu_{T/n}
\le n u^2(T/n)^2=u^2T^2/n.\]

Let \(n\to\infty\) to obtain zero variance at every \(T\). Write
\(\mu_t=\delta_{a(t)}\). The convolution law makes \(a\) additive, while
\(|a(t)|\le ut\) implies continuity at zero. Hence \(a(t)=bt\),
\(|b|\le u\). This proof even derives the needed continuity from the
support bound.

The proposition concerns stationary, spatially homogeneous, independent
increment composition on position alone. It leaves room for stochastic
bounded-speed trajectories with velocity memory. For example A01's
telegraph process is Markov in \((X,V)\); deleting \(V\) removes the
information needed for such position-only convolution. General
state-dependent Markov motion is also outside the stated convolution
hypothesis.

\hypertarget{consequence-for-the-cut-point-programme}{%
\subsection{4. Consequence for the cut-point
programme}\label{consequence-for-the-cut-point-programme}}

Finite propagation does not supply a positive microscopic midpoint
action coefficient in these models. It identifies a structural
requirement: a stochastic refinement law with strict speed must retain
correlations or additional state information, rather than resetting
independent position increments. The elastic construction makes this
visible at a single cut.

The next bounded task is to derive a finite-speed bridge on \((X,V)\),
including its conditional midpoint law and endpoint atoms, and compare
coarse and fine restrictions of the same experiment. Then separate a
finite observation-window action scale from a microscopic refinement
remainder. The universal positive scale remains a selection question
about the physical premises.

\hypertarget{proof-and-literature-record}{%
\subsection{Proof and literature
record}\label{proof-and-literature-record}}

The \href{../references/batches/B13.md}{B13 audit} reviews C030--C032 as
derived consequences. Barkai, \emph{Stable Equilibrium Based on Lévy
Statistics} (2003),
\href{https://arxiv.org/abs/cond-mat/0303255v1}{arXiv:cond-mat/0303255v1},
pp.~3--4, (2)--(3), supplies the elastic map. The equal-mass map is
well-defined even though the divided-density representation (5) excludes
that mass ratio. Cinque, \emph{A Note on the Conditional Probabilities
of the Telegraph Process} (2022),
\href{https://arxiv.org/abs/2202.01904v1}{arXiv:2202.01904v1}, pp.~1--2,
supplies the velocity-memory and endpoint-atom comparison. The midpoint
and convolution proofs above follow directly from their stated
hypotheses; neither selected source states those combined results. The
\href{../docs/batches/B13/B13-source-companion.md}{source companion}
records the bounded coverage. No priority claim follows from this audit.
\end{document}
